\chapter{Do gene à proteína}\label{ch:g11:gene-expression}

Em 1961 um bioquímico alimentou um extrato de bactérias sem \omterm{def:g10:cells-common-unit:cell}{células} com
um \omterm{def:g11:gene-expression:transcription}{RNA} artificial feito apenas da letra U, repetida. O extrato fabricou
uma \omterm{def:g11:gene-expression:protein}{proteína} feita de um único \omterm{def:g10:chemistry-of-life:families}{aminoácido}, a fenilalanina, repetido.
Com esse único experimento foi lida a primeira palavra do \omterm{def:g11:gene-expression:code}{código
genético}: UUU quer dizer fenilalanina. Em cinco anos as sessenta e
quatro palavras eram conhecidas, e mostraram ser as mesmas palavras
numa bactéria, num pé de trigo e num ser humano. Este capítulo segue o
caminho que vai da sequência de \omterm{def:g10:universal-dna:nucleotide}{nucleotídeos} de um \omterm{def:g10:universal-dna:gene}{gene} à sequência de
\omterm{def:g10:chemistry-of-life:families}{aminoácidos} de uma \omterm{def:g11:gene-expression:protein}{proteína} --- o caminho que toda característica do
\cref{ch:g11:mutations} percorre.

\section{Os genes fabricam proteínas}

\begin{proposition}[Um gene, uma proteína]\label{prop:g11:gene-expression:onegene}
Um \omterm{def:g10:universal-dna:gene}{gene} é expresso quando a \omterm{def:g10:cells-common-unit:cell}{célula} usa a sequência dele para fabricar a
\omterm{def:g11:gene-expression:protein}{proteína} que ele codifica. A maioria das características depende de
\omterm{def:g11:gene-expression:protein}{proteínas} --- \omterm{def:g10:cell-metabolism:metabolism}{enzimas}, fibras estruturais, transportadores, receptores
--- e uma \omterm{def:g11:mutations:mutation}{mutação} muda uma característica mudando a \omterm{def:g11:gene-expression:protein}{proteína} que o \omterm{def:g10:universal-dna:gene}{gene}
dela produz.
\end{proposition}
\begin{proof}[Evidências]
Beadle e Tatum (1941) irradiaram esporos de um bolor do pão e
recolheram mutantes que não conseguiam mais crescer em meio mínimo a
não ser que uma substância específica, como o \omterm{def:g10:chemistry-of-life:families}{aminoácido} arginina,
fosse acrescentada. Cada mutante não tinha uma \omterm{def:g10:cell-metabolism:metabolism}{enzima} da cadeia de
reações que fabrica essa substância; cada defeito era herdado como um
único \omterm{def:g10:universal-dna:gene}{gene}; e mutantes bloqueados em etapas diferentes da mesma cadeia
carregavam mutações em \omterm{def:g10:universal-dna:gene}{genes} diferentes. Um \omterm{def:g10:universal-dna:gene}{gene}, uma \omterm{def:g10:cell-metabolism:metabolism}{enzima} --- e,
como o trabalho posterior generalizou, um \omterm{def:g10:universal-dna:gene}{gene}, uma cadeia
polipeptídica.
\end{proof}

\begin{definition}[Proteína, sequência de aminoácidos]\label{def:g11:gene-expression:protein}
Uma \emph{proteína}\index{proteína} é uma cadeia de \omterm{def:g10:chemistry-of-life:families}{aminoácidos}, de
algumas dezenas a vários milhares de unidades, tirados de um conjunto
de vinte tipos. A \emph{sequência} --- a ordem dos \omterm{def:g10:chemistry-of-life:families}{aminoácidos} --- é o
que o \omterm{def:g10:universal-dna:gene}{gene} especifica; uma vez fabricada, a cadeia se dobra numa forma
tridimensional definida, determinada por essa sequência, e a forma
determina a função. Mude um \omterm{def:g10:chemistry-of-life:families}{aminoácido} e o dobramento, e a função,
podem mudar.
\end{definition}

\section{A transcrição: o gene copiado em RNA}

\begin{definition}[RNA mensageiro e transcrição]\label{def:g11:gene-expression:transcription}
O \emph{RNA}\index{RNA} (ácido ribonucleico) é uma cadeia de fita
simples de \omterm{def:g10:universal-dna:nucleotide}{nucleotídeos} como os do \omterm{def:g10:universal-dna:information}{DNA}, exceto que o açúcar dele é a
ribose e a base uracila (U) substitui a timina (T). A
\emph{transcrição}\index{transcrição} é a cópia da sequência de um \omterm{def:g10:universal-dna:gene}{gene}
numa molécula de \emph{RNA mensageiro}\index{RNA mensageiro} (mRNA): a
\omterm{def:g10:cell-metabolism:metabolism}{enzima} \emph{RNA polimerase} abre a \omterm{prop:g10:universal-dna:helix}{dupla-hélice} ao longo do \omterm{def:g10:universal-dna:gene}{gene}, lê
uma das fitas (a \emph{fita molde}) e monta o RNA complementar, U
diante de A, A diante de T, G diante de C e C diante de G. O mRNA tem a
sequência da outra fita do \omterm{def:g10:universal-dna:gene}{gene}, com U no lugar de T; ele sai do \omterm{def:g10:cells-common-unit:organelle}{núcleo}
para o \omterm{def:g10:cells-common-unit:cell}{citoplasma}.
\end{definition}

\begin{omfigure}
\begin{tikzpicture}[font=\small, >=Stealth, line cap=round]
  % DNA strands with bubble
  \draw[very thick, omDef] (0,0.5) -- (2.5,0.5) .. controls (3.2,0.5) and (3.4,1.3) .. (4.3,1.3) -- (6.7,1.3) .. controls (7.6,1.3) and (7.8,0.5) .. (8.5,0.5) -- (11,0.5);
  \draw[very thick, omDef] (0,-0.5) -- (2.5,-0.5) .. controls (3.2,-0.5) and (3.4,-1.3) .. (4.3,-1.3) -- (6.7,-1.3) .. controls (7.6,-1.3) and (7.8,-0.5) .. (8.5,-0.5) -- (11,-0.5);
  \foreach \x in {0.3,0.7,...,2.3} \draw[gray] (\x,0.5) -- (\x,-0.5);
  \foreach \x in {8.7,9.1,...,10.9} \draw[gray] (\x,0.5) -- (\x,-0.5);
  % template strand letters (bottom strand in bubble)
  \foreach \x/\l in {4.5/T, 4.9/A, 5.3/C, 5.7/G, 6.1/G, 6.5/A} \node[font=\scriptsize, omDef, anchor=north] at (\x,-1.32) {\l};
  % mRNA
  \draw[very thick, omThm] (4.4,-0.9) -- (6.6,-0.9) .. controls (7.4,-0.9) and (7.2,0.1) .. (8.0,0.6) .. controls (8.6,1.0) and (9.4,1.4) .. (10.6,1.7);
  \foreach \x/\l in {4.5/A, 4.9/U, 5.3/G, 5.7/C, 6.1/C, 6.5/U} \node[font=\scriptsize, omThm, anchor=south] at (\x,-0.85) {\l};
  \foreach \x in {4.5,4.9,...,6.5} \draw[gray, densely dotted] (\x,-0.95) -- (\x,-1.28);
  % polymerase
  \draw[thick, fill=omProp!30, rounded corners=8pt] (6.95,-1.6) rectangle (8.65,0.2);
  \node[align=center] at (7.8,-0.7) {RNA\\ polimerase};
  \draw[->, thick] (8.7,-1.9) -- (9.8,-1.9) node[right] {avança ao longo do gene};
  \node[anchor=east, omDef, align=right] at (-0.15,0) {DNA do\\ gene};
  \node[anchor=west, omThm, align=left] at (10.7,1.7) {mRNA, sendo\\ liberado};
  \node[anchor=north, omDef] at (5.5,-1.75) {fita molde};
\end{tikzpicture}

{\small A \omterm{def:g11:gene-expression:transcription}{transcrição}. A \omterm{def:g11:gene-expression:transcription}{RNA} polimerase abre a hélice, lê a fita molde
e constrói um \omterm{def:g11:gene-expression:transcription}{RNA mensageiro} complementar a ela (U no lugar de T).
Atrás da \omterm{def:g10:cell-metabolism:metabolism}{enzima} a hélice se fecha de novo; o \omterm{def:g11:gene-expression:transcription}{RNA} se desprende.}
\end{omfigure}

\begin{example}[Um transcrito]\label{ex:g11:gene-expression:transcript}
A fita molde \texttt{3'-TAC GGA CTT ATC-5'} dá o mRNA
\texttt{5'-AUG CCU GAA UAG-3'}. A outra fita do \omterm{def:g10:universal-dna:gene}{gene} diz
\texttt{5'-ATG CCT GAA TAG-3'}: o mRNA é a sequência dessa fita com U
no lugar de T, e é por isso que as sequências dos \omterm{def:g10:universal-dna:gene}{genes} são escritas
por convenção como essa fita, a \emph{fita codificadora}.
\end{example}

\section{O código genético}

\begin{definition}[Códon e código genético]\label{def:g11:gene-expression:code}
O mRNA é lido em grupos consecutivos de três \omterm{def:g10:universal-dna:nucleotide}{nucleotídeos}, os
\emph{códons}\index{códon}, a partir de um ponto de partida definido. O
\emph{código genético}\index{código genético} é a correspondência entre
os 64 códons possíveis e os 20 \omterm{def:g10:chemistry-of-life:families}{aminoácidos}: 61 códons especificam um
\omterm{def:g10:chemistry-of-life:families}{aminoácido}, três (UAA, UAG, UGA) são sinais de \emph{parada} que
encerram a cadeia, e AUG, que especifica a metionina, serve também de
códon de \emph{início}. O código é \emph{degenerado} --- a maioria dos
\omterm{def:g10:chemistry-of-life:families}{aminoácidos} tem vários códons --- e \emph{universal}: a mesma tabela
serve a todo organismo conhecido, com raras variantes menores.
\end{definition}

\begin{omfigure}
\begin{tabular}{c|cccc|c}
\multicolumn{1}{c}{} & \multicolumn{4}{c}{\small segunda letra} & \multicolumn{1}{c}{} \\
{\small primeira} & U & C & A & G & {\small terceira} \\ \hline
U & UUU Phe & UCU Ser & UAU Tyr & UGU Cys & U \\
 & UUC Phe & UCC Ser & UAC Tyr & UGC Cys & C \\
 & UUA Leu & UCA Ser & UAA \textbf{parada} & UGA \textbf{parada} & A \\
 & UUG Leu & UCG Ser & UAG \textbf{parada} & UGG Trp & G \\ \hline
C & CUU Leu & CCU Pro & CAU His & CGU Arg & U \\
 & CUC Leu & CCC Pro & CAC His & CGC Arg & C \\
 & CUA Leu & CCA Pro & CAA Gln & CGA Arg & A \\
 & CUG Leu & CCG Pro & CAG Gln & CGG Arg & G \\ \hline
A & AUU Ile & ACU Thr & AAU Asn & AGU Ser & U \\
 & AUC Ile & ACC Thr & AAC Asn & AGC Ser & C \\
 & AUA Ile & ACA Thr & AAA Lys & AGA Arg & A \\
 & AUG \textbf{Met} & ACG Thr & AAG Lys & AGG Arg & G \\ \hline
G & GUU Val & GCU Ala & GAU Asp & GGU Gly & U \\
 & GUC Val & GCC Ala & GAC Asp & GGC Gly & C \\
 & GUA Val & GCA Ala & GAA Glu & GGA Gly & A \\
 & GUG Val & GCG Ala & GAG Glu & GGG Gly & G \\
\end{tabular}

{\small O \omterm{def:g11:gene-expression:code}{código genético}, lido no mRNA. Cada \omterm{def:g11:gene-expression:code}{códon} é encontrado pela
primeira letra dele (linha do bloco), pela segunda (coluna) e pela
terceira (linha dentro do bloco). AUG é ao mesmo tempo a metionina e o
sinal de início; três \omterm{def:g11:gene-expression:code}{códons} são de parada.}
\end{omfigure}

\begin{proposition}[Como o código foi lido]\label{prop:g11:gene-expression:readcode}
O significado de cada \omterm{def:g11:gene-expression:code}{códon} foi estabelecido experimentalmente.
\end{proposition}
\begin{proof}[Evidências]
Nirenberg e Matthaei (1961) acrescentaram \omterm{def:g11:gene-expression:transcription}{RNA} sintético de um único
\omterm{def:g10:universal-dna:nucleotide}{nucleotídeo} repetido a um extrato sem \omterm{def:g10:cells-common-unit:cell}{células} que continha \omterm{def:g10:cells-common-unit:organelle}{ribossomos},
\omterm{def:g10:chemistry-of-life:families}{aminoácidos} e as \omterm{def:g10:cell-metabolism:metabolism}{enzimas} da síntese de \omterm{def:g11:gene-expression:protein}{proteínas}: o poli-U rendeu uma
cadeia de fenilalaninas, o poli-C uma cadeia de prolinas, o poli-A de
lisinas. \omterm{def:g11:gene-expression:transcription}{RNAs} sintéticos de pares e de trios repetidos e, mais tarde, a
ligação de trios isolados a \omterm{def:g10:cells-common-unit:organelle}{ribossomos} atribuíram os \omterm{def:g11:gene-expression:code}{códons} restantes
até 1966. O comprimento de três letras tinha sido deduzido da
aritmética (duas letras dão só 16 combinações, poucas demais para 20
\omterm{def:g10:chemistry-of-life:families}{aminoácidos}; três dão 64) e confirmado por mutações: inserir um ou dois
\omterm{def:g10:universal-dna:nucleotide}{nucleotídeos} num \omterm{def:g10:universal-dna:gene}{gene} destrói a \omterm{def:g11:gene-expression:protein}{proteína} dele, inserir três restaura
uma \omterm{def:g11:gene-expression:protein}{proteína} quase normal.
\end{proof}

\begin{method}[Traduzir uma sequência]\label{met:g11:gene-expression:translate}
\begin{enumerate}
  \item Escreva a fita codificadora (ou transcreva a fita molde):
        troque T por U para obter o mRNA.
  \item Encontre o primeiro AUG: ele é o início, e o quadro de leitura
        fica fixado a partir dele.
  \item Corte o mRNA em trios consecutivos a partir do AUG e leia cada
        um na tabela.
  \item Pare no primeiro \omterm{def:g11:gene-expression:code}{códon} de parada; os \omterm{def:g10:chemistry-of-life:families}{aminoácidos} listados, em
        ordem, são a sequência da \omterm{def:g11:gene-expression:protein}{proteína}.
  \item Para um mutante, repita e compare: mesma \omterm{def:g11:gene-expression:protein}{proteína}
        (silenciosa), um \omterm{def:g10:chemistry-of-life:families}{aminoácido} trocado (de sentido trocado),
        parada prematura (sem sentido), tudo mudado depois da \omterm{def:g11:mutations:mutation}{mutação}
        (deslocamento do quadro de leitura).
\end{enumerate}
\end{method}

\begin{example}[Quatro mutações, quatro desfechos]\label{ex:g11:gene-expression:outcomes}
Fita codificadora \texttt{ATG CCT GAA TTC TAG}: mRNA
\texttt{AUG CCU GAA UUC UAG}, \omterm{def:g11:gene-expression:protein}{proteína} Met--Pro--Glu--Phe.
\begin{itemize}
  \item \texttt{ATG CC\textbf{C} GAA TTC TAG}: CCC continua sendo Pro
        --- \emph{silenciosa}.
  \item \texttt{ATG CCT G\textbf{C}A TTC TAG}: GCA é Ala ---
        Met--Pro--Ala--Phe, \emph{de sentido trocado}.
  \item \texttt{ATG CCT \textbf{T}AA TTC TAG}: UAA é parada ---
        Met--Pro, \emph{sem sentido}, uma \omterm{def:g11:gene-expression:protein}{proteína} truncada.
  \item \texttt{ATG CC\textbf{A}T GAA TTC TAG}: um A inserido desloca
        o quadro: \texttt{AUG CCA UGA}\dots{} Met--Pro--parada ---
        \emph{deslocamento do quadro de leitura}.
\end{itemize}
\end{example}

\section{A tradução: a mensagem lida}

\begin{proposition}[A tradução]\label{prop:g11:gene-expression:translation}
A \emph{tradução}\index{tradução} acontece nos \emph{\omterm{def:g10:cells-common-unit:organelle}{ribossomos}}, no
\omterm{def:g10:cells-common-unit:cell}{citoplasma}. Um \omterm{def:g10:cells-common-unit:organelle}{ribossomo} se liga ao mRNA no \omterm{def:g11:gene-expression:code}{códon} de início dele e
avança ao longo dele, \omterm{def:g11:gene-expression:code}{códon} a \omterm{def:g11:gene-expression:code}{códon}. Cada \omterm{def:g11:gene-expression:code}{códon} é reconhecido por um
\emph{\omterm{def:g11:gene-expression:transcription}{RNA} transportador} (tRNA), um pequeno \omterm{def:g11:gene-expression:transcription}{RNA} que carrega, numa
extremidade, o \emph{anticódon} de três \omterm{def:g10:universal-dna:nucleotide}{nucleotídeos} complementar ao
\omterm{def:g11:gene-expression:code}{códon} e, na outra, o \omterm{def:g10:chemistry-of-life:families}{aminoácido} correspondente. O \omterm{def:g10:cells-common-unit:organelle}{ribossomo} une cada
\omterm{def:g10:chemistry-of-life:families}{aminoácido} trazido à cadeia em crescimento e libera o tRNA vazio; num
\omterm{def:g11:gene-expression:code}{códon} de parada ele libera a cadeia terminada, que se dobra. Vários
\omterm{def:g10:cells-common-unit:organelle}{ribossomos} leem um mesmo mRNA em sequência, e um mRNA rende centenas de
cópias da \omterm{def:g11:gene-expression:protein}{proteína} antes de ser degradado.
\end{proposition}
\admitted

\begin{omfigure}
\begin{tikzpicture}[font=\small, >=Stealth, line cap=round]
  % mRNA
  \draw[very thick, omThm] (0,0) -- (11,0);
  \foreach \x/\l in {0.4/A,0.8/U,1.2/G, 1.7/C,2.1/C,2.5/U, 3.0/G,3.4/A,3.8/A, 4.3/U,4.7/U,5.1/C, 5.6/G,6.0/G,6.4/U, 6.9/U,7.3/A,7.7/G}
    \node[font=\scriptsize, omThm, anchor=north] at (\x,-0.05) {\l};
  \node[anchor=east, omThm] at (-0.15,0) {mRNA};
  % ribosome (two lobes) over codons 3 and 4
  \draw[thick, fill=omProp!25, rounded corners=10pt] (2.7,0.15) rectangle (5.5,1.4);
  \draw[thick, fill=omProp!35, rounded corners=12pt] (2.5,1.4) .. controls (2.5,3.0) and (5.7,3.0) .. (5.7,1.4) -- cycle;
  \node at (4.1,2.1) {ribossomo};
  % tRNAs
  \draw[very thick, omMeth!80!black] (3.4,0.55) -- (3.4,1.35);
  \node[font=\scriptsize, omMeth!80!black, anchor=south] at (3.4,0.15) {CUU};
  \draw[very thick, omMeth!80!black] (4.7,0.55) -- (4.7,1.35);
  \node[font=\scriptsize, omMeth!80!black, anchor=south] at (4.7,0.15) {AAG};
  % amino acids chain
  \foreach \i/\x in {0/3.4, 1/4.7} \fill[omDef!70] (\x,1.5) circle (4pt);
  \foreach \x in {2.6, 1.9, 1.2} \fill[omDef!70] (\x,1.9) circle (4pt);
  \draw[thick, omDef!70] (1.2,1.9) -- (2.6,1.9) -- (3.4,1.5) -- (4.7,1.5);
  \node[anchor=east, omDef!80!black, align=right] at (0.9,1.9) {proteína\\ em crescimento};
  % incoming tRNA
  \draw[very thick, omMeth!80!black] (7.0,1.3) -- (7.0,2.1);
  \fill[omDef!70] (7.0,2.25) circle (4pt);
  \node[font=\scriptsize, omMeth!80!black, anchor=west] at (7.15,1.35) {CCA};
  \node[anchor=west, align=left, font=\scriptsize] at (7.3,1.7) {o tRNA seguinte, com o\\ aminoácido dele, chegando};
  \draw[->, thick] (7.0,1.1) -- (6.1,0.75);
  \draw[->, thick] (5.7,-0.9) -- (7.2,-0.9) node[right] {o ribossomo avança};
  \node[anchor=north, font=\scriptsize] at (4.05,-0.45) {códon};
  \draw (2.85,-0.35) -- (2.85,-0.45) -- (3.95,-0.45) -- (3.95,-0.35);
\end{tikzpicture}

{\small A \omterm{prop:g11:gene-expression:translation}{tradução}. O \omterm{def:g10:cells-common-unit:organelle}{ribossomo} abriga dois \omterm{def:g11:gene-expression:code}{códons}; um tRNA cujo
anticódon combina com cada \omterm{def:g11:gene-expression:code}{códon} traz o \omterm{def:g10:chemistry-of-life:families}{aminoácido} dele, a cadeia é
transferida para o recém-chegado, e o \omterm{def:g10:cells-common-unit:organelle}{ribossomo} avança um \omterm{def:g11:gene-expression:code}{códon}. Num
\omterm{def:g11:gene-expression:code}{códon} de parada a cadeia é liberada.}
\end{omfigure}

\begin{example}[Velocidade e produção]\label{ex:g11:gene-expression:speed}
Um \omterm{def:g10:cells-common-unit:organelle}{ribossomo} acrescenta de 5 a 20 \omterm{def:g10:chemistry-of-life:families}{aminoácidos} por segundo: uma \omterm{def:g11:gene-expression:protein}{proteína}
de 300 \omterm{def:g10:chemistry-of-life:families}{aminoácidos} leva menos de um minuto. Os \omterm{def:g10:cells-common-unit:organelle}{ribossomos} se seguem uns
aos outros num mRNA a cerca de 80 \omterm{def:g10:universal-dna:nucleotide}{nucleotídeos} de distância, de modo
que uma mensagem de 900 \omterm{def:g10:universal-dna:nucleotide}{nucleotídeos} carrega uma dúzia deles ao mesmo
tempo; ao longo da vida de algumas horas, um único mRNA produz mil
moléculas de \omterm{def:g11:gene-expression:protein}{proteína}. Uma \omterm{def:g10:cells-common-unit:cell}{célula} de \emph{E.~coli} fabrica umas
$\num{20000}$ moléculas de \omterm{def:g11:gene-expression:protein}{proteína} por segundo; a precursora de uma
hemácia, umas $\num{5e8}$ moléculas de hemoglobina ao longo da vida dela.
\end{example}

\section{Um gene, várias mensagens}

\begin{proposition}[Maturação da mensagem nos eucariontes]\label{prop:g11:gene-expression:splicing}
Nas \omterm{def:g10:cells-common-unit:prokaryote}{células eucarióticas} a sequência de um \omterm{def:g10:universal-dna:gene}{gene} é interrompida por
\emph{íntrons}\index{íntron}, segmentos que são transcritos mas depois
recortados do \omterm{def:g11:gene-expression:transcription}{RNA} antes de ele deixar o \omterm{def:g10:cells-common-unit:organelle}{núcleo}; os segmentos
conservados e unidos ponta a ponta, os \emph{éxons}, formam o mRNA
maduro. Muitos \omterm{def:g10:universal-dna:gene}{genes} podem ser recortados de mais de um jeito,
conservando conjuntos diferentes de éxons (o
\emph{splicing alternativo}), de modo que um \omterm{def:g10:universal-dna:gene}{gene} rende vários mRNAs
distintos e várias \omterm{def:g11:gene-expression:protein}{proteínas} aparentadas: os $\num{20000}$ \omterm{def:g10:universal-dna:gene}{genes}
humanos codificam bem mais de $\num{100000}$ \omterm{def:g11:gene-expression:protein}{proteínas}.
\end{proposition}
\admitted

\begin{omfigure}
\begin{tikzpicture}[font=\small, >=Stealth]
  \tikzset{ex/.style={draw, fill=omDef!40, minimum height=0.45cm, minimum width=1.0cm, inner sep=1pt},
           intr/.style={draw, fill=gray!20, minimum height=0.45cm, minimum width=0.7cm, inner sep=1pt}}
  \node[anchor=east] at (-0.2,0) {o gene (no DNA)};
  \node[ex] at (0.5,0) {1}; \node[intr] at (1.35,0) {}; \node[ex] at (2.2,0) {2}; \node[intr] at (3.05,0) {}; \node[ex] at (3.9,0) {3}; \node[intr] at (4.75,0) {}; \node[ex] at (5.6,0) {4};
  \node[anchor=west, font=\scriptsize] at (6.2,0) {éxons (em azul), íntrons (em cinza)};
  \draw[->, thick] (3.0,-0.35) -- (3.0,-0.95) node[midway, right] {transcrição};
  \node[anchor=east] at (-0.2,-1.3) {pré-mRNA};
  \node[ex, fill=omThm!30] at (0.5,-1.3) {1}; \node[intr] at (1.35,-1.3) {}; \node[ex, fill=omThm!30] at (2.2,-1.3) {2}; \node[intr] at (3.05,-1.3) {}; \node[ex, fill=omThm!30] at (3.9,-1.3) {3}; \node[intr] at (4.75,-1.3) {}; \node[ex, fill=omThm!30] at (5.6,-1.3) {4};
  \draw[->, thick] (2.0,-1.65) -- (1.0,-2.35) node[midway, left] {splicing};
  \draw[->, thick] (4.0,-1.65) -- (6.0,-2.35) node[midway, right] {splicing alternativo};
  \node[anchor=east] at (-0.6,-2.7) {mRNA A};
  \node[ex, fill=omThm!30] at (0.0,-2.7) {1}; \node[ex, fill=omThm!30] at (1.0,-2.7) {2}; \node[ex, fill=omThm!30] at (2.0,-2.7) {3}; \node[ex, fill=omThm!30] at (3.0,-2.7) {4};
  \node[anchor=west] at (3.6,-2.7) {proteína A};
  \node[anchor=east] at (5.4,-2.7) {};
  \node[ex, fill=omThm!30] at (6.0,-2.7) {1}; \node[ex, fill=omThm!30] at (7.0,-2.7) {2}; \node[ex, fill=omThm!30] at (8.0,-2.7) {4};
  \node[anchor=west] at (8.6,-2.7) {proteína B};
\end{tikzpicture}

{\small Um \omterm{def:g10:universal-dna:gene}{gene} eucariótico é transcrito inteiro, e depois os \omterm{prop:g11:gene-expression:splicing}{íntrons}
são removidos. Conservar todos os éxons, ou deixar um de fora, dá dois
mRNAs e duas \omterm{def:g11:gene-expression:protein}{proteínas} a partir do mesmo \omterm{def:g10:universal-dna:gene}{gene}.}
\end{omfigure}

\begin{remark}[O fluxo da informação]\label{rem:g11:gene-expression:flow}
O \omterm{def:g10:universal-dna:information}{DNA} é transcrito em \omterm{def:g11:gene-expression:transcription}{RNA}, o \omterm{def:g11:gene-expression:transcription}{RNA} é traduzido em \omterm{def:g11:gene-expression:protein}{proteína}, e a \omterm{def:g11:gene-expression:protein}{proteína}
faz a característica: a informação flui num sentido só. Uma \omterm{def:g11:gene-expression:protein}{proteína}
mudada nunca reescreve o \omterm{def:g10:universal-dna:gene}{gene}, e é por isso que uma vida inteira de
\omterm{def:g10:sport-and-health:training}{treinamento} ou de queimaduras de sol não é herdada, e por que só as
mutações do \omterm{def:g10:universal-dna:information}{DNA} (\cref{ch:g11:mutations}) mudam o que a geração
seguinte recebe. Os mesmos três passos correm em toda \omterm{def:g10:cells-common-unit:cell}{célula} de todo
organismo, no mesmo código: a universalidade do
\cref{ch:g10:universal-dna}, agora até a última palavra.
\end{remark}

\section{Exercícios}

\begin{exercise}[$\star$]\label{exo:g11:gene-expression:1}
Dê três diferenças entre o \omterm{def:g10:universal-dna:information}{DNA} e o \omterm{def:g11:gene-expression:transcription}{RNA}.
\end{exercise}

\begin{exercise}[$\star$]\label{exo:g11:gene-expression:2}
Transcreva a fita molde \texttt{3'-TAC CGA AAA ATT-5'} em mRNA.
\end{exercise}

\begin{exercise}[$\star$]\label{exo:g11:gene-expression:3}
Traduza o mRNA \texttt{AUG GGC AAA UGU UAA} com a tabela do código.
\end{exercise}

\begin{exercise}[$\star$]\label{exo:g11:gene-expression:4}
Quais são os papéis do \omterm{def:g10:cells-common-unit:organelle}{ribossomo} e dos \omterm{def:g11:gene-expression:transcription}{RNAs} transportadores na
\omterm{prop:g11:gene-expression:translation}{tradução}?
\end{exercise}

\begin{exercise}[$\star$]\label{exo:g11:gene-expression:5}
Por que o código precisa usar pelo menos três \omterm{def:g10:universal-dna:nucleotide}{nucleotídeos} por \omterm{def:g10:chemistry-of-life:families}{aminoácido}?
\end{exercise}

\begin{exercise}[$\star\star$]\label{exo:g11:gene-expression:6}
Uma \omterm{def:g11:gene-expression:protein}{proteína} tem 412 \omterm{def:g10:chemistry-of-life:families}{aminoácidos}. Qual é o comprimento mínimo da parte
codificadora do mRNA dela, \omterm{def:g11:gene-expression:code}{códon} de parada incluído, e do \omterm{def:g10:universal-dna:gene}{gene}
correspondente em pares de bases?
\end{exercise}

\begin{exercise}[$\star\star$]\label{exo:g11:gene-expression:7}
A fita codificadora \texttt{ATG AAA GGC TGG TAA} sofre \omterm{def:g11:mutations:mutation}{mutação} para
\texttt{ATG AAA GGC TGA TAA}. Dê as duas \omterm{def:g11:gene-expression:protein}{proteínas} e classifique a
\omterm{def:g11:mutations:mutation}{mutação}.
\end{exercise}

\begin{exercise}[$\star\star$]\label{exo:g11:gene-expression:8}
Mesma sequência de partida, mutada para \texttt{ATG AAG GGC TGG TAA} e
para \texttt{ATG AAA GGG CTG GTA A}. Classifique cada uma e dê as
\omterm{def:g11:gene-expression:protein}{proteínas}.
\end{exercise}

\begin{exercise}[$\star\star$]\label{exo:g11:gene-expression:9}
Explique por que uma substituição na terceira posição de um \omterm{def:g11:gene-expression:code}{códon} é
muitas vezes silenciosa, usando a tabela.
\end{exercise}

\begin{exercise}[$\star\star$]\label{exo:g11:gene-expression:10}
O \omterm{def:g11:gene-expression:transcription}{RNA} poli-UC (\texttt{UCUCUCUC}\dots) dá uma \omterm{def:g11:gene-expression:protein}{proteína} que alterna
serina e leucina. Mostre que isso é compatível com um código de três
letras e use o fato para atribuir dois \omterm{def:g11:gene-expression:code}{códons}.
\end{exercise}

\begin{exercise}[$\star\star$]\label{exo:g11:gene-expression:11}
Um \omterm{def:g10:universal-dna:gene}{gene} humano de \num{30000} pares de bases produz uma \omterm{def:g11:gene-expression:protein}{proteína} de 500
\omterm{def:g10:chemistry-of-life:families}{aminoácidos}. Explique a discrepância.
\end{exercise}

\begin{exercise}[$\star\star\star$]\label{exo:g11:gene-expression:12}
Um medicamento bloqueia os \omterm{def:g10:cells-common-unit:organelle}{ribossomos} bacterianos, mas não os humanos.
Explique por que ele pode ser um antibiótico, e o que a existência dele
implica sobre os \omterm{def:g10:cells-common-unit:organelle}{ribossomos} dos dois grupos.
\end{exercise}

\begin{exercise}[$\star\star\star$]\label{exo:g11:gene-expression:13}
Explique como o \omterm{def:g10:universal-dna:gene}{gene} da água-viva do \cref{ch:g10:universal-dna} pôde
ser lido por um camundongo, no vocabulário deste capítulo, e o que
aconteceria se o camundongo usasse outro código.
\end{exercise}

\begin{exercise}[$\star\star\star$]\label{exo:g11:gene-expression:14}
Uma \omterm{def:g11:mutations:mutation}{mutação} muda o tRNA cujo anticódon lê UAG, de modo que ele passa a
carregar um \omterm{def:g10:chemistry-of-life:families}{aminoácido} em vez de parar. Preveja o efeito dela sobre as
\omterm{def:g11:gene-expression:protein}{proteínas} da \omterm{def:g10:cells-common-unit:cell}{célula} em geral, e sobre um \omterm{def:g10:universal-dna:gene}{gene} mutante que carregue um
UAG prematuro.
\end{exercise}

\begin{exercise}[$\star\star\star$]\label{exo:g11:gene-expression:15}
Discuta por que um deslocamento do quadro de leitura perto do início de
um \omterm{def:g10:universal-dna:gene}{gene} é quase sempre pior que um perto do fim, e por que uma \omterm{def:g11:mutations:mutation}{mutação}
de sentido trocado pode ser qualquer coisa entre inofensiva e fatal.
\end{exercise}

\section{Problema: um gene lido de quatro maneiras}

\begin{problem}[{Problema de fim de semana --- um gene curto transcrito,
traduzido, mutado três vezes e cronometrado: do DNA aos aminoácidos, e
os sessenta segundos que uma proteína leva}]
\label{pb:g11:gene-expression:1}
A fita codificadora de um \omterm{def:g10:universal-dna:gene}{gene} curto diz:
\begin{center}
\texttt{ATG GCT TGG AAA CCC GAA TAC GGT CAT TTC TGA}
\end{center}

\medskip
\noindent\textbf{Parte I --- Ler o \omterm{def:g10:universal-dna:gene}{gene}.}
\begin{enumerate}
  \item Escreva a fita molde.
  \item Escreva o mRNA transcrito a partir do \omterm{def:g10:universal-dna:gene}{gene}.
  \item Traduza-o, \omterm{def:g11:gene-expression:code}{códon} a \omterm{def:g11:gene-expression:code}{códon}, e dê a sequência de \omterm{def:g10:chemistry-of-life:families}{aminoácidos} da
        \omterm{def:g11:gene-expression:protein}{proteína}.
  \item Quantos \omterm{def:g10:chemistry-of-life:families}{aminoácidos} a \omterm{def:g11:gene-expression:protein}{proteína} contém, e quantos \omterm{def:g10:universal-dna:nucleotide}{nucleotídeos}
        do mRNA são usados para especificá-los, parada incluída?
  \item Que \omterm{def:g11:gene-expression:code}{códon} começou a leitura, e o que fixa o quadro de leitura
        para todos os outros?
\end{enumerate}

\medskip
\noindent\textbf{Parte II --- Três mutantes.}
\begin{enumerate}[resume]
  \item Mutante 1: \texttt{ATG GCT TGG AAA CCG GAA TAC GGT CAT TTC TGA}.
        Dê a \omterm{def:g11:gene-expression:protein}{proteína} e classifique a \omterm{def:g11:mutations:mutation}{mutação}.
  \item Mutante 2: \texttt{ATG GCT TGG AAA CCC GAA TAC GGT CAT TTC TGA}
        com o nono \omterm{def:g11:gene-expression:code}{códon} \texttt{CAT} substituído por \texttt{CGT}.
        Dê a \omterm{def:g11:gene-expression:protein}{proteína} e classifique.
  \item Mutante 3: \texttt{ATG GCT TGA AAA CCC GAA TAC GGT CAT TTC TGA}.
        Dê a \omterm{def:g11:gene-expression:protein}{proteína} e classifique. Que substituição única a produziu?
  \item Mutante 4: um T é inserido depois do sexto \omterm{def:g10:universal-dna:nucleotide}{nucleotídeo}:
        \texttt{ATG GCT TTG GAA ACC CGA ATA CGG TCA TTT CTG A}. Dê a
        \omterm{def:g11:gene-expression:protein}{proteína} e classifique.
  \item Ordene os quatro mutantes do menos ao mais destrutivo para a
        \omterm{def:g11:gene-expression:protein}{proteína} e justifique.
\end{enumerate}

\medskip
\noindent\textbf{Parte III --- Por que o código é o que é.}
\begin{enumerate}[resume]
  \item Quantos \omterm{def:g11:gene-expression:code}{códons} especificam a leucina? E a serina? E o
        triptofano? Que \omterm{def:g10:chemistry-of-life:families}{aminoácido} tem mais chance de ser trocado por
        uma substituição aleatória no \omterm{def:g11:gene-expression:code}{códon} dele, e qual tem menos?
  \item Calcule a fração de todas as substituições na terceira posição
        de um \omterm{def:g11:gene-expression:code}{códon} de leucina \texttt{CUx} que são silenciosas.
  \item Explique por que um código de duas letras não funcionaria, e
        por que um código de quatro não é necessário.
  \item O poli-U dá poli-fenilalanina e o poli-A dá poli-lisina. Que
        entradas da tabela esses dois experimentos estabelecem?
  \item O código é o mesmo na bactéria e num ser humano. Enuncie a
        consequência para a \omterm{prop:g10:universal-dna:universal}{transgênese} e a consequência para o
        parentesco.
\end{enumerate}

\medskip
\noindent\textbf{Parte IV --- Cronometrar a fábrica.}
Um \omterm{def:g10:cells-common-unit:organelle}{ribossomo} acrescenta 10 \omterm{def:g10:chemistry-of-life:families}{aminoácidos} por segundo e os \omterm{def:g10:cells-common-unit:organelle}{ribossomos} se
seguem uns aos outros a 80 \omterm{def:g10:universal-dna:nucleotide}{nucleotídeos} de distância ao longo de um
mRNA; cada mRNA vive 2 horas.
\begin{enumerate}[resume]
  \item Quanto tempo um \omterm{def:g10:cells-common-unit:organelle}{ribossomo} leva para traduzir a \omterm{def:g11:gene-expression:protein}{proteína} deste
        \omterm{def:g10:universal-dna:gene}{gene}? E uma \omterm{def:g11:gene-expression:protein}{proteína} de 600 \omterm{def:g10:chemistry-of-life:families}{aminoácidos}?
  \item Quantos \omterm{def:g10:cells-common-unit:organelle}{ribossomos} conseguem ler ao mesmo tempo um mRNA de
        1800 \omterm{def:g10:universal-dna:nucleotide}{nucleotídeos}?
  \item Se um novo \omterm{def:g10:cells-common-unit:organelle}{ribossomo} começa a cada 8 segundos, quantas cópias
        da \omterm{def:g11:gene-expression:protein}{proteína} um mRNA rende ao longo da vida dele?
  \item Uma bactéria precisa de \num{2000} cópias de uma \omterm{def:g10:cell-metabolism:metabolism}{enzima} em 10
        minutos. Quantos mRNAs do \omterm{def:g10:universal-dna:gene}{gene}, transcritos ao mesmo tempo, são
        necessários?
  \item Enuncie o resultado: a \omterm{def:g11:gene-expression:protein}{proteína} codificada pelo \omterm{def:g10:universal-dna:gene}{gene}, a
        \omterm{def:g11:mutations:mutation}{mutação} entre as quatro que a aboliu, e o tempo de que um
        \omterm{def:g10:cells-common-unit:organelle}{ribossomo} precisa para construí-la.
\end{enumerate}
\end{problem}
